% Created 2024-12-11 Wed 16:59
% Intended LaTeX compiler: pdflatex
\documentclass[11pt]{article}
\PassOptionsToPackage{hyphens}{url}
\usepackage[utf8]{inputenc}
\usepackage[T1]{fontenc}
\usepackage{graphicx}
\usepackage{longtable}
\usepackage{wrapfig}
\usepackage{rotating}
\usepackage[normalem]{ulem}
\usepackage{amsmath}
\usepackage{amssymb}
\usepackage{capt-of}
\usepackage{hyperref}
\date{\today}
\title{Comments on "Development of an improvising machine in an ubimus context"}
\begin{document}

\maketitle
The system described appears to be interesting and useful for a variety of uses related to improvised music practices. I read the question on page 2, in the second paragraph, as the research question for the paper. This, however, does not appear to be correct. There is no further discussion of ensemble music practices following this section. Instead, the scope of the paper is far too wide for a short article like this. There are many small statements, such as the sentence following the question on page 2, that are not backed up by any data or even experience. (As a practitioner, I cannot see that in-person musical interactions have decreased post-covid.)

The definition of free improvisation is unsatisfying, and although several of the central sources are cited, there is a much more contemporary discussion of improvisation that could benefit the article's many references to free improvisation.

The aim, as described in the last paragraph of section 1 on page 3 is not followed up. The discussion in the paper does not properly contribute to the "broader discussion on the role of artificial agents in music" in the way that it is described here. It is not at all clear what the DIY perspective here is.

The description of the segmentation of musical parameters into lists on page 4 is far too general for the rest of the discussion to make sense. There are a number of questions not approached, the most important of which include time and segmentation. How are phrases identified and what constitutes a phrase? This can be radically different in different genres and styles of improvisation.

My suggestion for this paper would be to rewrite it into an overview of the field of automated systems for improvisation in the ubimus movement and relate these to the theory of free improvisation brought up in the paper, and then include a short discussion of the system developed by the authors or, to write a full description of the system and a study on when it is used in various contexts, how it is experienced by musicians, or how it is experienced by listeners. Either or all of those perspectives would be interesting.

Either way, the aim and scope of the paper needs to be properly stated.
\end{document}